\section{System Architecture}
\label{sec:architecture}

The pipeline has three layers. \Cref{fig:pipeline} shows them, the data that
moves between them, and the two points where state is durably written: the
\emph{datalake write} (header and body files, one of three directory layouts)
and the \emph{index write} (postings lists, one of three structures). Nothing
is cached in between: the index layer reads the bodies the datalake layer
just wrote, synchronously, in a single process.

\begin{figure}[htbp]
    \centering
    \begin{tikzpicture}[
        font=\small,
        layer/.style={draw, rounded corners, align=center, minimum height=2.2cm},
        state/.style={draw=black!60, dashed, rounded corners, align=center, inner sep=3pt},
        flow/.style={-{Stealth}, thick},
        label/.style={font=\footnotesize\itshape, text=black!70},
    ]
        % Layer 1: ingestion
        \node[layer, fill=blue!8, minimum width=4.6cm] (lake) {
            \textbf{Layer 1 --- Datalake}\\[2pt]
            download (sequential)\\
            header/body split\\
            layout: time / book / batch
        };
        % Layer 2: datamart
        \node[layer, fill=green!8, minimum width=4.6cm, right=1.6cm of lake] (mart) {
            \textbf{Layer 2 --- Datamart}\\[2pt]
            tokenize (shared rules)\\
            metadata (SQLite)\\
            inverted index builds
        };
        % Layer 3: serving
        \node[layer, fill=orange!10, minimum width=4.6cm, right=1.6cm of mart] (serve) {
            \textbf{Layer 3 --- Querying}\\[2pt]
            term lookup\\
            postings intersection\\
            document IDs returned
        };
        % Durable state
        \node[state, below=0.9cm of lake] (lakestate) {
            written state:\\
            \texttt{headers/\_, bodies/} (one layout)\\
            raw probe books (download test)
        };
        \node[state, below=0.9cm of mart] (martstate) {
            written state:\\
            \texttt{inverted\_index.json} $\mid$\\
            \texttt{inverted\_index/} (per-term files) $\mid$\\
            MongoDB collection \texttt{inverted\_index}
        };
        % Flows
        \draw[flow] (lake) -- node[label, above]{books (header, body)} (mart);
        \draw[flow] (mart) -- node[label, above]{loaded index, query terms} (serve);
        \draw[flow] (lake) -- (lakestate);
        \draw[flow] (mart) -- (martstate);
        \draw[flow, dashed] (lakestate.east) -- node[label, below]{re-read for indexing} (martstate.west);
    \end{tikzpicture}
    \caption{Stage 1 pipeline: three layers, the data that moves between them,
    and the durable state each layer writes. Downloads are sequential (one
    connection at a time); indexing reads the datalake from disk; querying
    only reads an already-built index.}
    \label{fig:pipeline}
\end{figure}

\paragraph{Layer 1 --- Datalake.}
For each book ID the engine downloads the raw Gutenberg text (or reads it
from the local corpus when offline), splits it at the official
\texttt{*** START/END OF THE PROJECT GUTENBERG EBOOK} markers, and stores the
licence header and the body as two files under one of three layouts:
\emph{time-based} (\texttt{year/month/day}), \emph{book-based}
(\texttt{book\_id/}), or \emph{batch/range-based}
(\texttt{batch\_start-batch\_end/}). The layout is a directory-structure
decision only; the stored bytes are identical.

\paragraph{Layer 2 --- Datamart.}
Bodies are tokenized with the shared rules of \cref{sec:design} and fed to
the inverted-index builders. Three structures are supported: a single JSON
file, a folder of per-term files, and a MongoDB collection. Metadata
(Title, Author, Language, paths) is extracted from the headers and persisted
in SQLite; the index builders read the corpus \emph{through} those metadata
records, which is the same path the production pipeline uses.

\paragraph{Layer 3 --- Querying.}
A query is a list of terms; a document matches when it contains \emph{every}
term. Each structure answers a term lookup with its postings list, and the
caller intersects those lists. This layer is exercised by the benchmark
workload; a user-facing search interface belongs to Stage 2 (see
\cref{sec:comment-matrix}, comments 32--33).

\paragraph{Synchronous by construction.}
The pipeline runs in one process, one book at a time: download, split, write,
tokenize, index. No step overlaps another, so every timed region in
\cref{sec:methodology} has an unambiguous start and stop. The only
concurrency anywhere in the measurement is the background thread that samples
the resident set size (\cref{tab:metrics}), and the MongoDB server process,
whose resources are explicitly excluded from client measurements
(\cref{sec:index-memory}).
